\chapter{Portfolio Construction II}\label{ch:rs:portfolio-construction-ii}

Move one stock's alpha by a basis point and a long-only mean--variance book of fifty names moves by 0.03 per cent of its capital on average, 0.33 per cent at
worst. Draw the forecast's noise again (an equally good forecast, with different errors) and the book turns over 24 per cent. Over eight years its
information about the stocks does not show: after costs, its Sharpe ratio of 0.57 is within one standard error of holding the market. Over days rather
than months the picture reverses: how a book moves matters more than where it points. With three true signals that decay at different speeds and trading
costs that grow with the square of the trade, re-optimising every day to the mean--variance portfolio nets a Sharpe ratio of $-0.51$; trading part of the way
each day towards Gârleanu and Pedersen's \omterm{def:rs:portfolio-construction-ii:aim}{aim portfolio} nets 4.06. This chapter builds \texttt{firm.allocation}: robust optimisation, Black--Litterman,
\omterm{def:rs:portfolio-construction-ii:erc}{risk parity} in two forms, and the \omterm{def:rs:portfolio-construction-ii:aim}{aim portfolio}.

\section{Robust optimisation}

\begin{definition}[Robust portfolio optimisation, uncertainty set]\label{def:rs:portfolio-construction-ii:robust}
\emph{Robust portfolio optimisation}\index{robust portfolio optimisation} maximises the worst case of the objective over an \emph{uncertainty
set}\index{uncertainty set} of inputs. For an ellipsoid of alphas $\{\alpha + \Omega^{1/2}u : \lVert u\rVert \le \kappa\}$ the worst-case expected return is
$\alpha^\top w - \kappa\sqrt{w^\top\Omega w}$, and the problem is a second-order cone programme (Goldfarb and Iyengar).
\end{definition}

The penalty $\kappa\sqrt{w^\top\Omega w}$ grows with the size of the book along the directions in which the alpha is least certain, which are the directions
the error maximiser of chapter~25 loved. The chapter's solver (\texttt{robust\_mv}) needs no cone solver; it works by majorisation: since $\sqrt x \le x/(2s) + s/2$
with equality at $x = s^2$, each step solves a quadratic programme with the square root replaced by a quadratic that touches it at the current book, and the
objective rises at every step (\cref{lst:rs:portfolio-construction-ii:robust}).

The chapter's first book holds the fifty largest names of chapter~25's universe, long only, fully invested, at most 10\% a name, rebuilt at the 95
month-ends with chapter~25's alpha and chapter~24's risk model ($\gamma = 20$). The robust version uses an ellipsoid of three standard errors of the alpha
across names. A basis point added to one stock's alpha moves the plain mean--variance book by 0.029\% of capital on average and 0.33\% at most; the robust
book by 0.028\% on average and 0.12\% at most. Robustness does not change how much a book listens to its forecast on average; it removes the largest jumps,
which happen where two names are nearly interchangeable to the optimiser.

\section{Views blended with equilibrium}

\begin{definition}[Implied equilibrium returns, Black--Litterman model]\label{def:rs:portfolio-construction-ii:bl}
The \emph{implied equilibrium returns}\index{implied equilibrium returns} of a market portfolio $w_m$ are $\Pi = \gamma\Sigma w_m$: the expected returns for
which $w_m$ is the mean--variance optimum. The \emph{Black--Litterman model}\index{Black--Litterman model} (Black and Litterman, 1992) treats $\Pi$ as a prior
with covariance $\tau\Sigma$ and views $Pr = q$ with uncertainty $\Omega$ as observations, and optimises on the posterior mean
$\mu = [(\tau\Sigma)^{-1} + P^\top\Omega^{-1}P]^{-1}[(\tau\Sigma)^{-1}\Pi + P^\top\Omega^{-1}q]$ with covariance $\Sigma + [(\tau\Sigma)^{-1} + P^\top\Omega^{-1}P]^{-1}$.
\end{definition}

Black--Litterman changes the question from ``what are the expected returns?'' to ``how far from the market should the forecast move me?'' Without views the
optimal book is the market; with a view held with certainty the posterior meets it exactly; the book in between is a blend. The chapter's version takes
capitalisation weights as $w_m$, each stock's alpha as a view on its return in excess of $\Pi$, $\tau = 0.05$ (the usual default) and view uncertainties equal to
$\tau$ times each stock's variance. Its sensitivity to one basis point is 0.031\% on average and 0.11\% at most, and its book stays near the market's: its
realised volatility is 19.6\%, the market's 19.5\%, where mean--variance's is 15.8\%.

\section{Risk parity and risk budgeting}

\begin{definition}[Risk contribution, risk budgeting, equal risk contribution portfolio, risk parity]\label{def:rs:portfolio-construction-ii:erc}
The \emph{risk contribution}\index{risk contribution} of position $i$ is $w_i(\Sigma w)_i/(w^\top\Sigma w)$, its Euler share of the variance (Book~6, chapter~20);
the shares sum to one. \emph{Risk budgeting}\index{risk budgeting} chooses long-only weights whose risk contributions equal given budgets; the \emph{equal risk
contribution portfolio}\index{equal risk contribution portfolio} gives every position the same share (Maillard, Roncalli and Teïletche). \emph{Risk
parity}\index{risk parity} is the practice of allocating across asset classes by equal risk contributions, usually with leverage.
\end{definition}

\begin{proposition}[Risk budgets by Newton's method]\label{prop:rs:portfolio-construction-ii:budget}
The function $f(y) = \frac12 y^\top\Sigma y - \sum_i b_i\ln y_i$ is strictly convex on $y > 0$ and its minimiser satisfies $y_i(\Sigma y)_i = b_i$; the weights
$w = y/\sum_i y_i$ therefore have \omterm{def:rs:portfolio-construction-ii:erc}{risk contributions} $b$.
\end{proposition}

\begin{proof}
The Hessian $\Sigma + \mathrm{diag}(b_i/y_i^2)$ is positive definite, and $f \to \infty$ at the boundary and at infinity, so there is one minimiser, where
$\Sigma y = b/y$ componentwise. \omterm{def:rs:portfolio-construction-ii:erc}{Risk contributions} are invariant to scaling $y$, and $\sum_i y_i(\Sigma y)_i = \sum_i b_i$.
\end{proof}

Risk-based books ignore the alpha, so they do not move when it does: their sensitivity is zero, by construction and without merit. Their virtue is elsewhere.
Maillard, Roncalli and Teïletche showed that the \omterm{def:rs:portfolio-construction-ii:erc}{equal risk contribution portfolio}'s volatility lies between the minimum-variance and the equally weighted
portfolios', and so it does here: 18.4\% against 19.3\% for equal weights. Asness, Frazzini and Pedersen explained why \omterm{def:rs:portfolio-construction-ii:erc}{risk parity} across asset classes has
earned more than its risk suggests: investors averse to leverage bid up risky assets, leaving safer ones with higher risk-adjusted returns for anyone who can
lever them.

\section{Hierarchical allocation}

\begin{definition}[Hierarchical risk parity]\label{def:rs:portfolio-construction-ii:hrp}
\emph{Hierarchical risk parity}\index{hierarchical risk parity} (López de Prado) orders the assets by single-linkage clustering of correlation distances
$\sqrt{(1-\rho_{ij})/2}$, and splits the capital recursively: each block of the ordered list is cut in two, and the halves receive capital in inverse proportion
to the variance of their inverse-variance portfolios.
\end{definition}

\omterm{def:rs:portfolio-construction-ii:hrp}{Hierarchical risk parity} never inverts the covariance matrix, so it cannot amplify its smallest eigenvalues; on a diagonal covariance it gives the
inverse-variance weights, and on a matrix of blocks it splits the capital between the blocks before it looks inside them. On the fifty names it holds at most
5.8\% of capital in one stock and realises 17.8\% volatility. The comparison of the eight rules over 95 months, after costs of ten basis points a trade:

\begin{center}
\small
\begin{tabular}{@{}lrrrrr@{}}
\toprule
rule & redraw turnover & Sharpe & after costs & volatility & turnover\\
\midrule
mean--variance, factor model & 23.9\% & 0.62 & 0.57 & 15.8\% & 33\%\\
mean--variance, sample covariance & 23.4\% & 0.82 & 0.77 & 15.9\% & 31\%\\
robust mean--variance & 22.4\% & 0.66 & 0.61 & 15.8\% & 31\%\\
Black--Litterman & 24.4\% & 0.75 & 0.72 & 19.6\% & 27\%\\
equal \omterm{def:rs:portfolio-construction-ii:erc}{risk contribution} & 0 & 0.66 & 0.65 & 18.4\% & 7\%\\
\omterm{def:rs:portfolio-construction-ii:hrp}{hierarchical risk parity} & 0 & 0.69 & 0.67 & 17.8\% & 14\%\\
equal weights & 0 & 0.67 & 0.66 & 19.3\% & 5\%\\
capitalisation weights & 0 & 0.66 & 0.65 & 19.5\% & 6\%\\
\bottomrule
\end{tabular}
\end{center}

The redraw turnover is the one-way turnover when the forecast's noise is drawn again; the last column is the monthly one-way turnover. A Sharpe ratio measured over 7.9 years has a standard error of about 0.4, and every row is within it of every other. The long-only books are dominated by the
market they hold; their alphas are a small tilt, and the forecast's noise, redrawn, moves the alpha-based books by a quarter of their capital. That is the
honest conclusion of a comparison that is often presented as a contest: the rules differ in stability, concentration and cost, not detectably in performance.

\begin{omfigure}
\begin{tikzpicture}
\begin{axis}[name=a,width=7cm,height=5.4cm,ybar,bar width=7pt,ylabel={\small turnover per redraw},xtick=data,symbolic x coords={MV,sample MV,robust,BL,ERC,
  HRP,equal,cap},x tick label style={rotate=45,anchor=east,font=\footnotesize},tick label style={font=\footnotesize},ymin=0,ymax=0.3,grid=major,
  grid style={gray!25},yticklabel style={/pgf/number format/fixed,/pgf/number format/precision=2},scaled y ticks=false,table/col sep=comma]
\addplot[fill=omDef!55,draw=omDef] table[x=label,y=redraw]{figdata/research/26-portfolio-construction-ii/methods.csv};
\end{axis}
\begin{axis}[at={(a.east)},anchor=west,xshift=1.2cm,width=7cm,height=5.4cm,ybar,bar width=7pt,ylabel={\small Sharpe ratio after costs},xtick=data,
  symbolic x coords={MV,sample MV,robust,BL,ERC,HRP,equal,cap},x tick label style={rotate=45,anchor=east,font=\footnotesize},tick label style={font=\footnotesize},
  ymin=0,ymax=1,grid=major,grid style={gray!25},table/col sep=comma]
\addplot[fill=gray!45,draw=gray] table[x=label,y=sr_net]{figdata/research/26-portfolio-construction-ii/methods.csv};
\end{axis}
\end{tikzpicture}

\omcaption{The eight long-only rules on fifty names: left, the one-way turnover when the forecast's noise is drawn again; right, the Sharpe ratio after costs
over 95 months (standard error about 0.4). Data: \texttt{rs\_allocation}.}\label{fig:rs:portfolio-construction-ii:methods}
\end{omfigure}

\section{Multi-period optimisation with costs}

\begin{definition}[Aim portfolio, multi-period optimisation]\label{def:rs:portfolio-construction-ii:aim}
\emph{Multi-period optimisation}\index{multi-period optimisation} chooses trades by planning a sequence of future books under forecasts of returns, risks
and costs, and executes the first (Boyd and co-authors). With quadratic costs $\frac\lambda2\Delta x^\top\Sigma\Delta x$ and signals whose predictive power
decays at rates $\phi_k$, Gârleanu and Pedersen show that the optimal book trades each period a fraction $a/\lambda$ of the way from the current book to the
\emph{aim portfolio}\index{aim portfolio} $(\gamma\Sigma)^{-1}\sum_k B_k f_k/(1 + \phi_k a/\gamma)$, which weights slow signals more than fast ones, with $a$ the
positive root of $a^2(1-\rho) + a(\gamma(1-\rho) + \lambda\rho) - \gamma\lambda(1-\rho) = 0$ and $\rho$ the discount rate.
\end{definition}

The chapter's second book holds the same fifty names long and short, daily, with the simulation's three planted expected-return components as signals: the
persistent drift (half-life 504 days, $\phi = 0.00137$ a day), the post-earnings drift (sixty days, $\phi = 0.0167$) and the one-day reversal ($\phi = 1$). They
are the true expected returns, so gross Sharpe ratios are high; the question is only what trading them costs. $\gamma = 72$ puts the daily Markowitz book at
10\% annual volatility. Four rules, at a cost level $\lambda = 100$:

\begin{center}
\small
\begin{tabular}{@{}lrrrr@{}}
\toprule
rule, $\lambda = 100$ & Sharpe ratio & after costs & costs a year & daily turnover\\
\midrule
re-optimise daily to Markowitz & 5.14 & $-0.51$ & 56.6\% & 2.11\\
\omterm{def:rs:portfolio-construction-ii:aim}{aim portfolio} (trade 0.56 of the gap) & 4.74 & 4.06 & 4.0\% & 0.56\\
partial adjustment to Markowitz (0.56) & 5.07 & 3.36 & 12.6\% & \\
rebalance monthly to the slow signals & 3.01 & 2.96 & 0.2\% & \\
\bottomrule
\end{tabular}
\end{center}

The daily Markowitz book chases the reversal, the largest and shortest of the three signals, and pays for it every day. Partial adjustment trades less but
still aims at the reversal. The \omterm{def:rs:portfolio-construction-ii:aim}{aim portfolio} aims mostly at the slow signals: with $a/\gamma = 0.78$, the reversal enters with weight $1/(1 + 0.78) = 0.56$, the
post-earnings drift with 0.987 and the persistent drift with 0.999, and the book moves 56\% of the way each day. At ten times the cost ($\lambda = 1\,000$)
the aim rule trades 0.23 of the gap and still nets 3.15; daily re-optimisation nets $-22.62$ with costs of 566\% a year, partial adjustment 1.08, the
monthly rebalance 2.55. At $\lambda = 10\,000$ only the aim rule survives, at 1.61 (\cref{fig:rs:portfolio-construction-ii:lambda}).

\begin{omfigure}
\begin{tikzpicture}
\begin{semilogxaxis}[width=11cm,height=6cm,xlabel={cost level $\lambda$},ylabel={Sharpe ratio after costs},tick label style={font=\small},label style={font=\small},
  xmin=8,xmax=12000,ymin=-3.6,ymax=5.2,grid=major,grid style={gray!25},legend columns=2,
  legend style={font=\footnotesize,at={(0.5,-0.25)},anchor=north,draw=none,fill=none,/tikz/every even column/.append style={column sep=8pt}},
  table/col sep=comma]
\addplot[black!60,thick,mark=o] table[x=lam,y=daily_Markowitz]{figdata/research/26-portfolio-construction-ii/lambda.csv};
\addplot[omDef,very thick,mark=*] table[x=lam,y=aim_portfolio]{figdata/research/26-portfolio-construction-ii/lambda.csv};
\addplot[omDef!55,thick,mark=square*] table[x=lam,y=partial_adjustment]{figdata/research/26-portfolio-construction-ii/lambda.csv};
\addplot[omThm,thick,dashed,mark=triangle*] table[x=lam,y=monthly_Markowitz]{figdata/research/26-portfolio-construction-ii/lambda.csv};
\draw[black!50,densely dotted] (axis cs:8,0) -- (axis cs:12000,0);
\legend{daily Markowitz,aim portfolio,partial adjustment,monthly Markowitz}
\end{semilogxaxis}
\end{tikzpicture}

\omcaption{Sharpe ratio after quadratic costs against the cost level, for four trading rules on the three planted signals; values below $-3.5$ are drawn at
$-3.5$. Data: \texttt{rs\_allocation.multiperiod}.}\label{fig:rs:portfolio-construction-ii:lambda}
\end{omfigure}

\section{Tutorial: the one-basis-point flip}\label{tut:rs:portfolio-construction-ii:tut}

\begin{tutorial}
\textbf{Goal.} Build the fifty-name book by eight rules, measure their stability and realised performance, then trade the three planted signals daily under
four rules and a range of costs. \textbf{End state:} the two tables, \cref{fig:rs:portfolio-construction-ii:methods,fig:rs:portfolio-construction-ii:lambda}.
\begin{tutsteps}
\item \textbf{Robust mean--variance} by majorisation.
\omcode{code/firm/allocation/firm_allocation.py}{48}{59}{A sequence of quadratic programmes for the robust objective.}\label{lst:rs:portfolio-construction-ii:robust}
\item \textbf{Risk budgets} by Newton's method.
\omcode{code/firm/allocation/firm_allocation.py}{85}{101}{Weights with given risk contributions.}\label{lst:rs:portfolio-construction-ii:budget}
\item \textbf{The \omterm{def:rs:portfolio-construction-ii:aim}{aim portfolio} and the trade rate.}
\omcode{code/firm/allocation/firm_allocation.py}{157}{170}{Gârleanu and Pedersen's rule.}\label{lst:rs:portfolio-construction-ii:gp}
\item \textbf{Run} \texttt{rs\_allocation.stability}, \texttt{redraw\_turnover}, \texttt{backtest} for each rule, \texttt{multiperiod} for each rule and cost, and
\texttt{fig\_allocation.py}.
\end{tutsteps}
\textbf{What to change next.} Replace the true signals by noisy forecasts and see the aim rule's advantage grow (it averages the noise, as chapter~25's turnover limit
did); give the equal-risk-contribution book leverage to the mean--variance book's volatility and compare.
\end{tutorial}

\section{Build: the allocation toolkit}\label{bld:rs:portfolio-construction-ii:allocation}

\begin{build}
\bfield{Purpose} The alternatives to plain mean--variance that the firm uses, each with a test of its defining property, and a trading rule for signals of
different speeds.
\bfield{Interface} \texttt{mean\_variance(alpha, Sigma, gamma, lo, hi, budget)}, \texttt{robust\_mv(alpha, Sigma, Omega, kappa, gamma, lo, hi, budget)},
\texttt{implied\_returns(Sigma, w\_mkt, gamma)}, \texttt{black\_litterman(pi, Sigma, P, q, Omega, tau)}, \texttt{risk\_contributions(w, Sigma)},
\texttt{risk\_budget(Sigma, b)}, \texttt{erc(Sigma)}, \texttt{hrp(Sigma)}, \texttt{gp\_trade\_rate(gamma, lam, rho)}, \texttt{gp\_aim(Sigma, signals, phis, gamma, a)}.
\bfield{Rules} Every rule is judged on stability to a redrawn forecast and on cost-adjusted performance with its standard error; a signal's decay rate is
estimated (chapter~13) before it is traded; fast signals are traded only through a rule that knows they are fast.
\bfield{Acceptance tests} \texttt{code/firm/allocation/tests/}: mean--variance against its closed form; the robust objective above the plain book's, and equal
to it at $\kappa = 0$; implied returns reproducing the market; a certain view met and a worthless one ignored; \omterm{def:rs:portfolio-construction-ii:erc}{risk contributions} equal to their budgets; HRP
equal to inverse variance on a diagonal matrix and splitting blocks evenly; the trade rate's limits and the aim's down-weighting.
\bfield{Stretch} The aim rule with linear costs (no-trade regions); robust optimisation over the covariance too; hierarchical risk budgets.
\end{build}

\begin{omsources}
\item F.~Black and R.~Litterman, ``Global portfolio optimization'', \emph{Financial Analysts Journal} 48(5), 1992; PyPortfolioOpt, \texttt{black\_litterman}
(documented implementation, after He and Litterman).
\item D.~Goldfarb and G.~Iyengar, ``Robust portfolio selection problems'', \emph{Mathematics of Operations Research} 28(1), 2003.
\item S.~Maillard, T.~Roncalli and J.~Teïletche, ``The properties of equally weighted risk contribution portfolios'', \emph{Journal of Portfolio Management}
36(4), 2010.
\item M.~López de Prado, ``Building diversified portfolios that outperform out of sample'', \emph{Journal of Portfolio Management} 42(4), 2016.
\item N.~Gârleanu and L.~H.~Pedersen, ``Dynamic trading with predictable returns and transaction costs'', \emph{Journal of Finance} 68(6), 2013.
\item S.~Boyd et al., ``Multi-period trading via convex optimization'', \emph{Foundations and Trends in Optimization} 3(1), 2017.
\item C.~S.~Asness, A.~Frazzini and L.~H.~Pedersen, ``Leverage aversion and risk parity'', \emph{Financial Analysts Journal} 68(1), 2012.
\end{omsources}

\section{Exercises}

\begin{exercise}[$\star$]\label{exo:rs:portfolio-construction-ii:1}
Two uncorrelated assets have volatilities of 20\% and 10\%. Give the equal-weight portfolio's volatility, and the weights of the \omterm{def:rs:portfolio-construction-ii:erc}{equal risk contribution portfolio}.
\end{exercise}

\begin{exercise}[$\star$]\label{exo:rs:portfolio-construction-ii:2}
Show that with no views the Black--Litterman book is the market portfolio.
\end{exercise}

\begin{exercise}[$\star$]\label{exo:rs:portfolio-construction-ii:3}
With $a/\gamma = 0.78$, give the \omterm{def:rs:portfolio-construction-ii:aim}{aim portfolio}'s weights on signals decaying at $\phi = 1$, $1/60$ and $0.00137$ a day.
\end{exercise}

\begin{exercise}[$\star\star$]\label{exo:rs:portfolio-construction-ii:4}
Why does robust optimisation cut the largest response to a one-basis-point change but not the average one?
\end{exercise}

\begin{exercise}[$\star\star$]\label{exo:rs:portfolio-construction-ii:5}
Why is the \omterm{def:rs:portfolio-construction-ii:erc}{equal risk contribution portfolio}'s volatility between the minimum-variance and the equally weighted portfolios'?
\end{exercise}

\begin{exercise}[$\star\star$]\label{exo:rs:portfolio-construction-ii:6}
Why does partial adjustment to the Markowitz portfolio do worse after costs than the aim rule trading at the same rate?
\end{exercise}

\begin{exercise}[$\star\star\star$]\label{exo:rs:portfolio-construction-ii:7}
\emph{Coding.} Run \texttt{rs\_allocation.multiperiod} for the four rules at $\lambda = 1\,000$ and $10\,000$, and explain which rule survives and why.
\end{exercise}

\begin{exercise}[$\star\star\star$]\label{exo:rs:portfolio-construction-ii:8}
\emph{Find the flaw.} ``\omterm{def:rs:portfolio-construction-ii:hrp}{Hierarchical risk parity} beat mean--variance by 0.10 of Sharpe ratio after costs over eight years, so we will switch.''
\end{exercise}

\section{Problem: The One-Basis-Point Flip}

\begin{problem}[{Weekend problem --- stable books and good trades}]\label{pb:rs:portfolio-construction-ii:1}
The chapter's fifty-name books and daily multi-period rules.

\medskip\noindent\textbf{Part I --- Stability.}
\begin{enumerate}
\item How much does a basis point on one alpha move each alpha-based book, on average and at most?
\item How much does a redrawn forecast move them?
\item Why do the risk-based books not move, and why is that no merit?
\item What does the robust ellipsoid remove, and what not?
\end{enumerate}

\medskip\noindent\textbf{Part II --- Equilibrium and risk.}
\begin{enumerate}[resume]
\item Write the Black--Litterman posterior mean and explain $\tau$.
\item What are the equal \omterm{def:rs:portfolio-construction-ii:erc}{risk contribution} and HRP books' volatilities and largest weights?
\item How does HRP avoid inverting the covariance matrix?
\item Why can \omterm{def:rs:portfolio-construction-ii:erc}{risk parity} across asset classes earn more than its risk suggests?
\end{enumerate}

\medskip\noindent\textbf{Part III --- Performance.}
\begin{enumerate}[resume]
\item Give the eight rules' Sharpe ratios after costs and their standard error.
\item What can and cannot be concluded from the comparison?
\end{enumerate}

\medskip\noindent\textbf{Part IV --- Trading through time.}
\begin{enumerate}[resume]
\item What are the three signals and their decay rates?
\item Why does daily re-optimisation lose after costs?
\item What does the \omterm{def:rs:portfolio-construction-ii:aim}{aim portfolio} weight, and how fast does it trade at $\lambda = 100$?
\item State the \emph{named result}: turnover per basis point of alpha perturbation for each method, and the net Sharpe ratio of the \omterm{def:rs:portfolio-construction-ii:aim}{aim portfolio} against daily
re-optimisation.
\item How do the rules behave as costs grow tenfold and a hundredfold?
\item Why does the monthly rebalance do well at low costs and badly at high ones?
\item What would noisy forecasts change?
\item Which of the chapter's rules would you run for a daily signal book?
\item How would you estimate $\lambda$?
\item In one sentence: what matters more at a daily horizon, the book or the path to it?
\end{enumerate}
\end{problem}

\section{Interview questions}

\begin{interviewq}[$\star$ \iqroles{researcher}]\label{iq:rs:portfolio-construction-ii:1}
What problem does Black--Litterman solve, and how?
\end{interviewq}

\begin{interviewq}[$\star\star$ \iqroles{researcher, risk}]\label{iq:rs:portfolio-construction-ii:2}
Explain \omterm{def:rs:portfolio-construction-ii:erc}{risk parity}. What does it assume, and what does it need to deliver equity-like returns?
\end{interviewq}

\begin{interviewq}[$\star\star$ \iqroles{researcher}]\label{iq:rs:portfolio-construction-ii:3}
How does robust optimisation change a mean--variance portfolio? How would you choose the size of the \omterm{def:rs:portfolio-construction-ii:robust}{uncertainty set}?
\end{interviewq}

\begin{interviewq}[$\star\star$ \iqroles{researcher}]\label{iq:rs:portfolio-construction-ii:4}
Describe \omterm{def:rs:portfolio-construction-ii:hrp}{hierarchical risk parity} and when you would use it.
\end{interviewq}

\begin{interviewq}[$\star\star$ \iqroles{researcher, trader}]\label{iq:rs:portfolio-construction-ii:5}
You have a fast signal and a slow signal and quadratic trading costs. How should you trade them?
\end{interviewq}

\begin{interviewq}[$\star\star\star$ \iqroles{researcher}]\label{iq:rs:portfolio-construction-ii:6}
Derive the weights of the \omterm{def:rs:portfolio-construction-ii:erc}{equal risk contribution portfolio} for two assets, and show that with uncorrelated assets they are inverse-volatility weights.
\end{interviewq}
